% GENERATED FILE. Edit canonical lessons, then run npm run generate:books.
% canonical-source-hash: fnv1a64:6851ed31c2491a31
% canonical-lessons: GU-C387-bahetar, GU-C387-uttam, GU-C387-cadiyatu, GU-C387-saras, GU-C387-kathalvu

\chapter{Better and Worse}
\label{ch:gu-better-and-worse}
\hlchaptermodality{387}

\begin{chapteropening}
\textbf{By the end of this chapter:} \emph{I can say whether things will get better or worse.}

Complete the last lesson of chapter 387: I can say whether things will get better or worse.
\end{chapteropening}

\section[\texorpdfstring{\gu{બહેતર} (bahetar)}{બહેતર (bahetar)}]{\gu{બહેતર} (bahetar) — better}

\label{lesson:GU-C387-bahetar}



\begin{quote}
Before the new one: say the Gujarati for subject to, then the Gujarati for to give in.
\end{quote}

\subsection*{You'll want to know: \gu{બહેતર}}
\textbf{\gu{બહેતર}} — \emph{bahetar} — \textquotedblleft{}better\textquotedblright{}.

\begin{quote}
\gu{જો} \gu{તમે} \gu{વહેલા} \gu{સૂશો} \gu{તો} \gu{તમારી} \gu{તબિયત} \gu{બહેતર} \gu{થશે}.
\end{quote}

\begin{quote}
\emph{jo tame vahelā sūśo to tamārī tabiyat bahetar thaśe.}
\end{quote}

\begin{quote}
If you go to bed early, your health will improve.
\end{quote}

\begin{cousinweb}
\textbf{\gu{બહેતર}} (\emph{bahetar}) is usually traced to Persian \emph{behtar}, \textquotedblleft{}better\textquotedblright{}.
\end{cousinweb}

\begin{grammarlens}[title={What you've built}]
One more word for saying what happens if.
\end{grammarlens}

\subsection*{Your turn}
\begin{itemize}
\raggedright
  \item \emph{Say it:} \emph{bahetar}
  \item \emph{Say it:} \emph{bahetar}, once more
  \item \emph{Say it:} \emph{jo tame vahelā sūśo to tamārī tabiyat bahetar thaśe.} — If you go to bed early, your health will improve.
  \item \emph{Recall:} say the Gujarati for subject to, then the Gujarati for to give in, then say \emph{bahetar} again
\end{itemize}

\begin{tcolorbox}[breakable,colback=teal!4,colframe=teal!35!black,title={Before you move on}]
What does \gu{બહેતર} mean? (\textquotedblleft{}Better\textquotedblright{}.) Say it once more.
\end{tcolorbox}
\section[\texorpdfstring{\gu{ઉત્તમ} (uttam)}{ઉત્તમ (uttam)}]{\gu{ઉત્તમ} (uttam) — excellent}

\label{lesson:GU-C387-uttam}



\begin{quote}
Before the new one: say the Gujarati for to give in, then the Gujarati for better.
\end{quote}

\subsection*{You'll want to know: \gu{ઉત્તમ}}
\textbf{\gu{ઉત્તમ}} — \emph{uttam} — \textquotedblleft{}excellent\textquotedblright{}.

\begin{quote}
\gu{જો} \gu{તમે} \gu{તાજું} \gu{શાક} \gu{વાપરશો} \gu{તો} \gu{દાળ} \gu{ઉત્તમ} \gu{બનશે}.
\end{quote}

\begin{quote}
\emph{jo tame tājũ śāk vāparśo to dāḷ uttam banśe.}
\end{quote}

\begin{quote}
If you use fresh vegetables, the dal will turn out excellent.
\end{quote}

\begin{cousinweb}
\textbf{\gu{ઉત્તમ}} (\emph{uttam}) is a tatsama, borrowed from Sanskrit \emph{uttama}, \textquotedblleft{}uppermost, highest, best\textquotedblright{} (Monier-Williams).
\end{cousinweb}

\begin{grammarlens}[title={What you've built}]
One more word for saying what happens if.
\end{grammarlens}

\subsection*{Your turn}
\begin{itemize}
\raggedright
  \item \emph{Say it:} \emph{uttam}
  \item \emph{Say it:} \emph{uttam}, once more
  \item \emph{Say it:} \emph{jo tame tājũ śāk vāparśo to dāḷ uttam banśe.} — If you use fresh vegetables, the dal will turn out excellent.
  \item \emph{Recall:} say the Gujarati for to give in, then the Gujarati for better, then say \emph{uttam} again
\end{itemize}

\begin{tcolorbox}[breakable,colback=teal!4,colframe=teal!35!black,title={Before you move on}]
What does \gu{ઉત્તમ} mean? (\textquotedblleft{}Excellent\textquotedblright{}.) Say it once more.
\end{tcolorbox}
\section[\texorpdfstring{\gu{ચડિયાતું} (caḍiyātũ)}{ચડિયાતું (caḍiyātũ)}]{\gu{ચડિયાતું} (caḍiyātũ) — superior}

\label{lesson:GU-C387-cadiyatu}



\begin{quote}
Before the new one: say the Gujarati for better, then the Gujarati for excellent.
\end{quote}

\subsection*{You'll want to know: \gu{ચડિયાતું}}
\textbf{\gu{ચડિયાતું}} — \emph{caḍiyātũ} — \textquotedblleft{}superior\textquotedblright{}.

\begin{quote}
\gu{જો} \gu{તમે} \gu{રોજ} \gu{કસરત} \gu{કરશો} \gu{તો} \gu{તમે} \gu{બીજા} \gu{ખેલાડીઓ} \gu{કરતાં} \gu{ચડિયાતા} \gu{થશો}.
\end{quote}

\begin{quote}
\emph{jo tame roj kasrat karśo to tame bījā khelāḍīo kartā̃ caḍiyātā thaśo.}
\end{quote}

\begin{quote}
If you exercise every day, you will be better than the other players.
\end{quote}

\begin{cousinweb}
\textbf{\gu{ચડિયાતું}} (\emph{caḍiyātũ}) is built on \emph{caḍvũ}, \textquotedblleft{}to climb, to rise\textquotedblright{}, the same verb as \textbf{\gu{ચઢવું}} (\emph{caḍhvũ}), from a reconstructed \emph{caḍhati}, \textquotedblleft{}rises\textquotedblright{} (CDIAL 4578): one who has climbed higher.
\end{cousinweb}

\begin{grammarlens}[title={What you've built}]
One more word for saying what happens if.
\end{grammarlens}

\subsection*{Your turn}
\begin{itemize}
\raggedright
  \item \emph{Say it:} \emph{caḍiyātũ}
  \item \emph{Say it:} \emph{caḍiyātũ}, once more
  \item \emph{Say it:} \emph{jo tame roj kasrat karśo to tame bījā khelāḍīo kartā̃ caḍiyātā thaśo.} — If you exercise every day, you will be better than the other players.
  \item \emph{Recall:} say the Gujarati for better, then the Gujarati for excellent, then say \emph{caḍiyātũ} again
\end{itemize}

\begin{tcolorbox}[breakable,colback=teal!4,colframe=teal!35!black,title={Before you move on}]
What does \gu{ચડિયાતું} mean? (\textquotedblleft{}Superior\textquotedblright{}.) Say it once more.
\end{tcolorbox}
\section[\texorpdfstring{\gu{સરસ} (saras)}{સરસ (saras)}]{\gu{સરસ} (saras) — lovely; fine}

\label{lesson:GU-C387-saras}



\begin{quote}
Before the new one: say the Gujarati for excellent, then the Gujarati for superior.
\end{quote}

\subsection*{You'll want to know: \gu{સરસ}}
\textbf{\gu{સરસ}} — \emph{saras} — \textquotedblleft{}lovely; fine\textquotedblright{}.

\begin{quote}
\gu{જો} \gu{કાલે} \gu{તડકો} \gu{હશે} \gu{તો} \gu{ઉજાણીનો} \gu{સરસ} \gu{દિવસ} \gu{હશે}.
\end{quote}

\begin{quote}
\emph{jo kāle taḍko haśe to ujāṇīno saras divas haśe.}
\end{quote}

\begin{quote}
If it is sunny tomorrow, it will be a lovely day for a picnic.
\end{quote}

\begin{cousinweb}
\textbf{\gu{સરસ}} (\emph{saras}) is inherited from Sanskrit \emph{sárasa}, \textquotedblleft{}juicy; moist, beautiful\textquotedblright{}; CDIAL lists the Gujarati word, \textquotedblleft{}juicy, excellent\textquotedblright{} (CDIAL 13255).
\end{cousinweb}

\begin{grammarlens}[title={What you've built}]
One more word for saying what happens if.
\end{grammarlens}

\subsection*{Your turn}
\begin{itemize}
\raggedright
  \item \emph{Say it:} \emph{saras}
  \item \emph{Say it:} \emph{saras}, once more
  \item \emph{Say it:} \emph{jo kāle taḍko haśe to ujāṇīno saras divas haśe.} — If it is sunny tomorrow, it will be a lovely day for a picnic.
  \item \emph{Recall:} say the Gujarati for excellent, then the Gujarati for superior, then say \emph{saras} again
\end{itemize}

\begin{tcolorbox}[breakable,colback=teal!4,colframe=teal!35!black,title={Before you move on}]
What does \gu{સરસ} mean? (\textquotedblleft{}Lovely; fine\textquotedblright{}.) Say it once more.
\end{tcolorbox}
\section[\texorpdfstring{\gu{કથળવું} (kathaḷvũ)}{કથળવું (kathaḷvũ)}]{\gu{કથળવું} (kathaḷvũ) — to deteriorate}

\label{lesson:GU-C387-kathalvu}



\begin{quote}
Before the new one: say the Gujarati for superior, then the Gujarati for lovely.
\end{quote}

\subsection*{You'll want to know: \gu{કથળવું}}
\textbf{\gu{કથળવું}} — \emph{kathaḷvũ} — \textquotedblleft{}to deteriorate\textquotedblright{}.

\begin{quote}
\gu{જો} \gu{ખોટ} \gu{વધશે} \gu{તો} \gu{ધંધો} \gu{કથળશે}.
\end{quote}

\begin{quote}
\emph{jo khoṭ vadhśe to dhandho kathaḷśe.}
\end{quote}

\begin{quote}
If the losses grow, the business will go downhill.
\end{quote}

Past (it …): \gu{તે} \gu{ધંધો} \gu{કથળ્યો}, \emph{te dhandho kathaḷyo}.

\begin{cousinweb}
\textbf{\gu{કથળવું}} (\emph{kathaḷvũ}) CDIAL lists the Gujarati verb, \textquotedblleft{}to be put out of order, to be spoilt\textquotedblright{}, under a reconstructed \emph{kusthala}, \textquotedblleft{}a bad place\textquotedblright{} (CDIAL 3383).
\end{cousinweb}

\begin{grammarlens}[title={What you've built}]
One more word for saying what happens if.
\end{grammarlens}

\subsection*{Your turn}
\begin{itemize}
\raggedright
  \item \emph{Say it:} \emph{kathaḷvũ}
  \item \emph{Say it:} \emph{kathaḷvũ}, once more
  \item \emph{Say it:} \emph{jo khoṭ vadhśe to dhandho kathaḷśe.} — If the losses grow, the business will go downhill.
  \item \emph{Recall:} say the Gujarati for superior, then the Gujarati for lovely, then say \emph{kathaḷvũ} again
\end{itemize}

\begin{tcolorbox}[breakable,colback=teal!4,colframe=teal!35!black,title={Before you move on}]
What does \gu{કથળવું} mean? (\textquotedblleft{}To deteriorate\textquotedblright{}.) Say it once more.
\end{tcolorbox}
